% Generated draft: requires mathematical and source review.
\documentclass[11pt]{article}
\usepackage[margin=25mm]{geometry}
\usepackage{fontspec}
\setmainfont{Latin Modern Roman}
\usepackage{amsmath,amssymb,mathtools,unicode-math}
\setmathfont{Latin Modern Math}
\usepackage{enumitem,needspace,xurl,hyperref}
\setlength{\parindent}{0pt}
\setlength{\parskip}{5pt}
\setlength{\emergencystretch}{3em}
\newcommand{\N}{\mathbb{N}}
\newcommand{\Z}{\mathbb{Z}}
\newcommand{\Q}{\mathbb{Q}}
\newcommand{\R}{\mathbb{R}}
\newcommand{\C}{\mathbb{C}}
\newcommand{\F}{\mathbb{F}}
\newcommand{\A}{\mathbb{A}}
\newcommand{\PP}{\mathbb{P}}
\newcommand{\E}{\mathbb{E}}
\newcommand{\eps}{\varepsilon}
\newcommand{\dd}{\,\mathrm{d}}
\begin{document}
\section*{UnsolvedMath ID: 2301021}
\textbf{Source problem code:} catalogue entry
\subsection{Definitions and mathematical statement}
Let \(f\) be meromorphic on \(\C\), and let \(T(r,f)\) be its Nevanlinna characteristic: the mean of \(\log^+|f|\) on \(|z|=r\) plus the integrated pole-counting function. Whenever the eventual denominator is positive, define extended-real quantities
\[
\alpha_f=\limsup_{r\to\infty}\frac{T(r,f)}{T(r,f')},\qquad
\beta_f=\liminf_{r\to\infty}\frac{T(r,f)}{T(r,f')}.
\]
The source asks whether there exists an admissible \(f\) with \(\beta_f>1\), and whether there exists one with \(\beta_f=\infty\). The admissible class is not precise enough as transcribed: it says only ``non-constant in the plane.'' This includes polynomials. For example, the characteristic asymptotics of \(f(z)=z^2\) and \(f'(z)=2z\) give ratio tending to \(2\), so the reported uncertainty about \(\beta_f>1\) cannot concern that literal class. Affine functions can also have \(T(r,f')=0\), leaving the quotient undefined. A missing restriction such as transcendentality is plausible, but is not silently imposed here. After fixing the class and denominator convention, the two displayed existence questions are the intended logical targets.

The source records the sharp known inequalities
\[
\alpha_f\ge\tfrac12\quad\text{for meromorphic }f,
\qquad \alpha_f\ge1\quad\text{for entire }f,
\]
and says \(\beta_f\) is finite for finite-order functions. Its class qualifications need the same review just noted. It also reports examples of entire functions of every order \(0\le\rho\le\infty\) with \(\alpha_f=\infty\), and, for every \(K>0\) and \(\rho>0\), an entire function of order at most \(\rho\) such that
\[
\frac{T(r,f)}{T(r,f')}>K
\]
on a set \(E\subseteq[1,\infty)\) with positive lower logarithmic density,
\[
\liminf_{R\to\infty}\frac1{\log R}\int_{E\cap[1,R]}\frac{\dd r}{r}>0.
\]
Order here means \(\limsup_{r\to\infty}\log T(r,f)/\log r\). The supplied OPEN-TRIAGE summary describes the question as open in the 2018 edition, with partial constructions and no recent resolution located; this does not resolve the transcription's missing function-class hypothesis.
\subsection{Short English statement}
Can the lower limiting ratio of a function's characteristic to its derivative's characteristic exceed one, or be infinite? The transcribed class includes polynomial examples and zero denominators, so an intended restriction, likely excluding polynomials, must be checked.
\subsection{Sources}
\begin{itemize}
\item Local UnsolvedMath registry, permanent ID \(2301021\), exact source code \texttt{catalogue entry}, ``Research Problems in Function Theory --- Problem 1.21.''
\item Hayman, \emph{Research Problems in Function Theory} (2018), Problem 1.21; \url{https://arxiv.org/abs/1809.07200}. Source-TeX extraction; recorded access 2026-07-29. The statement refers to further Hayman sources without complete citations.
\item Supplied OPEN-TRIAGE summary: open as of the 2018 edition, partial constructions, no recent resolution located; no separate current-status citation supplied.
\end{itemize}
\end{document}
